\subsection{Mitigating Hallucinations}

A notable finding is that FP2MP does not achieve the highest knowledge integration score (3.77 $\pm$ 0.55 vs. ReAct’s 4.67 $\pm$ 0.47).
We attribute this to a design trade-off: ReAct explicitly externalizes retrieved text, while FP2MP delegates deterministic operations to executable code and tool use (CodeSpatialAgent, BlocksNetAgent).
The case study shows 79\% of tool calls go to BlocksNet, producing verifiable numeric evidence (e.g., 45.59 min average accessibility).
This makes FP2MP more robust to spatial hallucinations but less "verbally rich" in knowledge display.

This design choice also enables extensibility as FP2MP does not require re-implementing domain logic.
Any existing urban analytics tool can be integrated as a new agent.
The orchestrator will automatically route relevant subqueries, and the blackboard will curate its outputs alongside other agents.
BlocksNetAgent serves as a concrete demonstration of this pattern.